\scanpage{71}
\margin{- 140 -}
\margin{IX}

où par définition $({}_n A_K)^f=({}_n A^o)_K^f$ désigne le r.ème composant $(n=\ell^{r+1})$ du pro-groupe de Barsotti-Tate $T_\ell(A_K)^f$. La première des relations (5.5.8) nous fournit alors un isomorphisme canonique
\[
(11.6.7)\qquad {}_n\Psi_K\simeq\underline M_{nK}
\qquad\text{où }\underline M_n\overset{\mathrm{dfn}}=
\underline M\otimes_{\mathbb Z}(\mathbb Z/n\mathbb Z).
\]
En fait, on a :

\uline{Lemme} 11.6.8. \uline{Le schéma en groupes} $ {}_n\Psi$ (11.6.6) \uline{sur} $S$ \uline{est étale sur} $S$, \uline{et l'isomorphisme} (11.6.7) \uline{se prolonge de façon unique en un morphisme}
\[
(11.6.8.1)\qquad {}_n\Psi\lhook\joinrel\longrightarrow\underline M_n
\]
\uline{qui est une immersion ouverte}.

Pour voir que $ {}_n\Psi$ est étale, il reste à vérifier en vertu de (11.6.7) qu'il l'est en les points de sa fibre spéciale, ou encore que cette dernière est étale, ce qui est trivial, la fibre en question étant
$ {}_n A_o/{}_n A_o^o\subset A_o/A_o^o=\Phi_o$.
Il s'ensuit que $ {}_n\Psi$ est un schéma normal, et comme $\underline M_n$ est fini sur $S$, l'existence d'un prolongement (11.6.8.1) en résulte aussitôt. C'est un morphisme de schémas en groupes étales sur $S$, donc son noyau est un groupe étale sur $S$, et comme $ {}_n\Psi$ donc ce noyau est séparé, et sa fibre générique est le groupe unité par (11.6.7), c'est le groupe unité (3.1.3). Donc (11.6.8.1) est un monomorphisme étale, donc une immersion ouverte (EGA IV 17.9.1), ce qui prouve le lemme.

Posons maintenant
\[
(11.6.9)\qquad {}_n\Phi=({}_n A)^f/({}_n A^o)^f\subset{}_n\Psi,
\]
de sorte que $ {}_n\Phi$ est un sous-schéma fini ouvert de $ {}_n\Psi$, puisque $({}_n A)^f$

\margin{458}
\margin{458}
\margin{70}
